\chapter{Neural Networks, CNNs, and Sequence Models}
\label{ch:deep}

\chapterguide%
  {Logistic regression from \chapref{ch:classification}, gradient descent from \chapref{ch:learning}, overfitting and validation from \chapref{ch:training}.}
  {explain layers, training, CNNs, and recurrent models; validate networks against simpler baselines.}

Neural networks combine weighted layers with nonlinear activations. CNNs suit image grids; recurrent models suit sequences. Keep the same split--fit--validate--test workflow.

\section{From logistic regression to a neural network}

A \term{layer} computes weighted sums, applies nonlinear \term{activations}, and passes outputs onward. Use a logistic/softmax output for classification or a linear output for regression.

\begin{mlfigure}
\begin{tikzpicture}[
  n/.style={circle,draw=MLBlue,fill=MLPaleBlue,minimum size=6.5mm,inner sep=0pt,font=\scriptsize\color{MLNavy}},
  h/.style={circle,draw=MLTeal,fill=MLPaleTeal,minimum size=6.5mm,inner sep=0pt,font=\scriptsize\color{MLNavy}},
  o/.style={circle,draw=MLOrange,fill=MLPaleOrange,minimum size=6.5mm,inner sep=0pt,font=\scriptsize\color{MLNavy}},
  e/.style={draw=MLMuted!55,thin}
]
\foreach \i/\y in {1/1.6,2/0.8,3/0,4/-0.8} \node[n] (x\i) at (0,\y) {$x_\i$};
\foreach \j/\y in {1/1.6,2/0.8,3/0,4/-0.8,5/-1.6} \node[h] (h\j) at (2.6,\y) {};
\foreach \k/\y in {1/0.8,2/0,3/-0.8} \node[h] (g\k) at (5.2,\y) {};
\node[o] (out) at (7.8,0) {$\hat{y}$};
\foreach \i in {1,...,4} \foreach \j in {1,...,5} \draw[e] (x\i) -- (h\j);
\foreach \j in {1,...,5} \foreach \k in {1,...,3} \draw[e] (h\j) -- (g\k);
\foreach \k in {1,...,3} \draw[e] (g\k) -- (out);
\node[font=\scriptsize\color{MLMuted}] at (0,-1.6) {inputs};
\node[font=\scriptsize\color{MLMuted},align=center] at (2.6,-2.4) {hidden layer 1\\(5 weighted sums, each\\passed through ReLU)};
\node[font=\scriptsize\color{MLMuted}] at (5.2,-1.6) {hidden layer 2};
\node[font=\scriptsize\color{MLMuted}] at (7.8,-1.0) {output};
\end{tikzpicture}
\end{mlfigure}

\term{ReLU}, $\max(0,z)$, retains positive inputs and zeros negative ones. Without nonlinear activations, stacked linear layers collapse to one linear map.

\begin{workedexample}
Take two inputs $x=(2,1)$, a hidden layer with two units and weights $w_1=(1,-1)$, $w_2=(0.5,0.5)$, no biases, and ReLU. The two sums are $2-1=1$ and $1+0.5=1.5$; ReLU leaves both as they are. An output unit with weights $(1,-2)$ then gives $1-3=-2$; through a sigmoid that is a probability of about $0.12$. Change $x$ to $(1,2)$: the first sum becomes $-1$, ReLU sets it to $0$, and the output becomes $0-3=-3$. The zero is the non-linearity at work: the network can behave differently on the two sides of the line $x_1=x_2$, which no single linear model can do.
\end{workedexample}

\begin{intuition}
Hidden units create nonlinear features; later layers combine them into richer decision boundaries.
\end{intuition}

\section{How a network is trained}

Train by gradient-based optimization (\secref{sec:gradient-descent}):

\begin{mlsteps}
  \item \term{Forward pass.} Predict a minibatch and compute classification cross-entropy or regression squared loss.
  \item \term{Backward pass.} Backpropagation applies the chain rule to compute every weight's gradient.
  \item \term{Update.} An optimizer uses those gradients; the learning rate controls step size. The labs use Adam.
  \item \term{Repeat.} One pass through training batches is an \term{epoch}. Record validation performance each epoch.
\end{mlsteps}

\begin{mltable}
\begin{tabularx}{\linewidth}{@{}>{\bfseries\leavevmode\color{MLNavy}\ttfamily}p{0.2\linewidth}p{0.18\linewidth}X@{}}
\toprule
\rowcolor{MLPaleBlue!92}
\normalfont\bfseries\leavevmode\color{MLNavy}Setting & \normalfont\bfseries\leavevmode\color{MLNavy}Typical value & \normalfont\bfseries\leavevmode\color{MLNavy}What raising it does \\
\midrule
learning rate & 0.001--0.01 & Bigger steps: faster at first, unstable if too big \\
batch size & 32--128 & Smoother gradients, fewer updates per epoch \\
epochs & 10--50 & More passes; watch validation, not training \\
hidden units & 32--256 & More capacity; more risk of memorising \\
dropout & 0.1--0.5 & Randomly silences units during training; a regularizer \\
\bottomrule
\end{tabularx}
\end{mltable}

Monitor overfitting (\chapref{ch:training}). \term{Early stopping} retains the best validation weights; \term{dropout} randomly disables units during training.

\begin{keyidea}
Extra capacity makes validation and a one-time final test more important, not less.
\end{keyidea}

\begin{pitfall}
Scale inputs before training and save any training-fitted scaling statistics.
\end{pitfall}

\section{Convolutional networks for images}

A \term{convolutional layer} slides learned local \term{filters} across an image. Each weighted response forms part of a \term{feature map}.

\begin{mlfigure}
\begin{tikzpicture}[font=\scriptsize]
% image grid
\foreach \i in {0,...,5} \foreach \j in {0,...,5} \draw[MLMuted!50] (\j*0.5,-\i*0.5) rectangle ++(0.5,-0.5);
\fill[MLPaleBlue,opacity=0.7] (0.5,-0.5) rectangle (2.0,-2.0);
\draw[MLBlue,very thick] (0.5,-0.5) rectangle (2.0,-2.0);
\node[font=\small\bfseries\color{MLNavy}] at (1.5,0.35) {image $8\times8$};
\node[color=MLBlue] at (1.25,-1.25) {$3\times3$};
\draw[-{Stealth[length=1.6mm]},MLBlue,thick] (2.05,-1.25) -- (2.55,-1.25);
\draw[-{Stealth[length=1.6mm]},MLBlue,thick] (1.25,-2.05) -- (1.25,-2.55);
% filter
\begin{scope}[xshift=4.2cm]
\foreach \i in {0,...,2} \foreach \j in {0,...,2} \draw[MLMuted!60] (\j*0.5,-\i*0.5) rectangle ++(0.5,-0.5);
\node at (0.25,-0.25) {$-1$}; \node at (0.75,-0.25) {$0$}; \node at (1.25,-0.25) {$1$};
\node at (0.25,-0.75) {$-1$}; \node at (0.75,-0.75) {$0$}; \node at (1.25,-0.75) {$1$};
\node at (0.25,-1.25) {$-1$}; \node at (0.75,-1.25) {$0$}; \node at (1.25,-1.25) {$1$};
\node[font=\small\bfseries\color{MLNavy}] at (0.75,0.35) {one filter};
\node[color=MLMuted,align=center] at (0.75,-2.05) {this one responds\\to vertical edges};
\end{scope}
% feature map
\begin{scope}[xshift=7.4cm]
\foreach \i in {0,...,5} \foreach \j in {0,...,5} \draw[MLMuted!50] (\j*0.5,-\i*0.5) rectangle ++(0.5,-0.5);
\fill[MLTeal!70] (1.0,-0.5) rectangle (1.5,-3.0);
\fill[MLTeal!35] (1.5,-0.5) rectangle (2.0,-3.0);
\node[font=\small\bfseries\color{MLNavy}] at (1.5,0.35) {feature map};
\node[color=MLMuted,align=center] at (1.5,-3.4) {bright where the\\pattern occurs};
\end{scope}
\draw[-{Stealth[length=2mm]},MLTeal,thick] (3.1,-1.5) -- (4.1,-1.5);
\draw[-{Stealth[length=2mm]},MLTeal,thick] (6.0,-1.5) -- (7.3,-1.5);
\end{tikzpicture}
\end{mlfigure}

Filters share weights across locations. \term{Max pooling} retains maxima within blocks, reducing spatial size. A small CNN repeats convolution--ReLU--pooling, then flattens maps into a dense classifier.

\begin{workedexample}
Lab~7 uses 16 then 32 filters on $8\times8$ digits. Two pooling stages leave $2\times2$ maps: $32\times2\times2=128$ values feed the ten-class classifier. The network trains on a CPU.
\end{workedexample}

\begin{intuition}
Training learns filters that reduce loss, rather than relying on hand-specified edges or shapes.
\end{intuition}

\section{Sequence models: RNNs and LSTMs}

A \term{recurrent network} processes ordered inputs using a \term{hidden state}. Each step combines the new input with the previous state; the resulting state supports prediction.

\begin{mlfigure}
\begin{tikzpicture}[
  cell/.style={draw=MLTeal,fill=MLPaleTeal,rounded corners=1.5mm,minimum width=1.5cm,minimum height=0.9cm,font=\scriptsize\bfseries\color{MLNavy}},
  inp/.style={draw=MLBlue,fill=MLPaleBlue,rounded corners=1.2mm,minimum width=1.2cm,minimum height=0.6cm,font=\scriptsize\color{MLNavy}},
  outp/.style={draw=MLOrange,fill=MLPaleOrange,rounded corners=1.2mm,minimum width=1.5cm,minimum height=0.6cm,font=\scriptsize\color{MLNavy}},
  arr/.style={-{Stealth[length=1.8mm]},thick,draw=MLTeal}
]
\foreach \t/\x in {1/0,2/2.4,3/4.8,24/8.4} {
  \node[cell] (c\t) at (\x,0) {LSTM};
  \node[inp] (i\t) at (\x,-1.4) {$x_{\t}$};
  \draw[arr] (i\t) -- (c\t);
}
\node[font=\Large\color{MLMuted}] at (6.6,0) {$\cdots$};
\draw[arr] (c1) -- node[above,font=\scriptsize\color{MLMuted}] {state} (c2);
\draw[arr] (c2) -- node[above,font=\scriptsize\color{MLMuted}] {state} (c3);
\draw[arr] (c3) -- (6.1,0);
\draw[arr] (7.1,0) -- (c24);
\node[outp] (o) at (8.4,1.5) {$\hat{y}_{25}$};
\draw[arr] (c24) -- (o);
\node[font=\scriptsize\color{MLMuted},align=center] at (4.2,-2.3) {24 hourly demand values go in, one step at a time; the same cell (same weights) is used at every step;\\the final state predicts the next hour};
\end{tikzpicture}
\end{mlfigure}

An \term{LSTM} adds a memory cell and learned forget, input, and output gates. These control what to retain, write, and expose, helping preserve longer-term dependencies.

Lab~8 predicts hour $t$ from the previous 24 demand values. Sliding windows form supervised rows; split chronologically to prevent future leakage.

\begin{pitfall}
Compare against last-value and same-hour-yesterday forecasts. Keep the network only if its validated gain justifies its cost.
\end{pitfall}

\section{Practical notes}

\begin{itemize}
  \item \term{Set the seed} (\texttt{torch.manual\_seed}) and expect small run-to-run differences anyway; training is stochastic.
  \item \term{Scale inputs} using training data only; save the fitted scaling.
  \item \term{Save weights} with \texttt{torch.save(model.state\_dict(), path)}. Include the architecture and model card needed to reload them.
  \item \term{Windows OpenMP conflicts:} duplicate runtimes may trigger ``OMP: Error \#15''. Check the environment setup before rerunning.
  \item \term{A CPU suffices} for these labs; larger workloads may benefit from a GPU.
  \item \term{Keep classical baselines.} Compare neural models with boosting on tables and simpler predictors on images or sequences.
\end{itemize}

\begin{modelcard}{Neural Networks - Quick Guide}
\begin{tabularx}{\linewidth}{@{}>{\bfseries\leavevmode\color{MLNavy}}p{0.22\linewidth}X@{}}
Best when & Inputs are images, sequences, or text; there is plenty of data; useful features are hard to write by hand. \cr
Start with & A small dense network as the reference, then a CNN for grids or an LSTM for sequences; always beside a classical baseline. \cr
Preprocessing & Scale inputs to a small range using training data; window a time series; split chronologically when rows overlap in time. \cr
Main controls & Learning rate, epochs (with early stopping), batch size, layer sizes, dropout. \cr
Watch for & Memorising the training set, unscaled inputs, run-to-run variation, and a baseline that was never beaten. \cr
\end{tabularx}
\end{modelcard}

\begin{libnote}
Labs~7--8 use \lib{PyTorch} tensors, \texttt{nn.Sequential}, and explicit training loops. \lib{Keras} instead offers a higher-level \texttt{model.fit} interface.
\end{libnote}

\labsection{7}{Deep learning I: a first neural network and a CNN for images}{sec:lab07-cnn}

\begin{labproject}{Train a dense network and a convolutional network on the handwritten digits}
\textbf{Goal.} Build tensors, train dense/CNN models, compare Lab~4's baseline, inspect errors and filters, and save weights with a card.

\medskip\noindent\textbf{Setup.}\par
\begin{lstlisting}[style=mlpython]
import json
import os
import zipfile
from pathlib import Path

# Anaconda's NumPy and pip's PyTorch each ship an OpenMP runtime on Windows; let them coexist.
os.environ.setdefault("KMP_DUPLICATE_LIB_OK", "TRUE")

import matplotlib.pyplot as plt
import numpy as np
import pandas as pd
import torch
from sklearn.metrics import ConfusionMatrixDisplay, accuracy_score
from torch import nn

torch.manual_seed(42)
PROJECT_ROOT = Path.cwd().parent
DATA = PROJECT_ROOT / "all_datasets"
MODELS = PROJECT_ROOT / "models"
MODELS.mkdir(exist_ok=True)
print("PyTorch", torch.__version__, "| device: cpu")
\end{lstlisting}

\medskip\noindent\textbf{Part A --- Images as tensors.}\par

Reshape 64-pixel digits to $1\times8\times8$ and divide brightnesses by 16. CNN tensors use (batch, channels, height, width). Reserve one-quarter of training rows for validation and the official test file for final evaluation.

\begin{lstlisting}[style=mlpython]
columns = [f"pixel_{i}" for i in range(64)] + ["label"]
with zipfile.ZipFile(DATA / "optical_recognition_handwritten_digits.zip") as archive:
    members = {Path(name).name: name for name in archive.namelist()}
    train = pd.read_csv(archive.open(members["optdigits.tra"]), header=None, names=columns)
    test = pd.read_csv(archive.open(members["optdigits.tes"]), header=None, names=columns)

def to_tensors(frame):
    X = torch.tensor(frame.drop(columns="label").to_numpy(), dtype=torch.float32)
    X = X.reshape(-1, 1, 8, 8) / 16.0            # (rows, channels, height, width), values in [0, 1]
    y = torch.tensor(frame["label"].to_numpy(), dtype=torch.long)
    return X, y

X_dev, y_dev = to_tensors(train)
X_test, y_test = to_tensors(test)
perm = torch.randperm(len(X_dev), generator=torch.Generator().manual_seed(42))
n_val = len(X_dev) // 4
X_val, y_val = X_dev[perm[:n_val]], y_dev[perm[:n_val]]
X_fit, y_fit = X_dev[perm[n_val:]], y_dev[perm[n_val:]]
print("fit:", X_fit.shape, " validation:", X_val.shape, " test:", X_test.shape)

fig, axes = plt.subplots(1, 10, figsize=(10, 1.4))
for ax, image, label in zip(axes, X_fit[:10], y_fit[:10]):
    ax.imshow(image[0], cmap="gray_r")
    ax.set_title(int(label), fontsize=9)
    ax.axis("off")
plt.suptitle("Ten training digits as the network sees them", y=1.05)
plt.show()
\end{lstlisting}

\medskip\noindent\textbf{Part B --- A first neural network and the training loop.}\par

The dense model maps 64 pixels through 128 ReLU units to ten scores. Logistic regression is the no-hidden-layer baseline.

\begin{lstlisting}[style=mlpython]
def make_mlp():
    return nn.Sequential(
        nn.Flatten(),                 # 1 x 8 x 8  ->  64
        nn.Linear(64, 128), nn.ReLU(),
        nn.Linear(128, 10),           # ten class scores
    )

mlp = make_mlp()
print(mlp)
print("trainable weights:", sum(p.numel() for p in mlp.parameters()))
\end{lstlisting}

Each epoch shuffles rows into batches of 32. Compute scores, cross-entropy, gradients via \texttt{backward()}, and optimizer updates; record training/validation accuracy.

\begin{lstlisting}[style=mlpython]
def accuracy(model, X, y):
    model.eval()
    with torch.no_grad():
        return (model(X).argmax(dim=1) == y).float().mean().item()

def train_classifier(model, X, y, X_v, y_v, epochs=20, lr=3e-3, batch_size=32):
    optimizer = torch.optim.Adam(model.parameters(), lr=lr)
    loss_fn = nn.CrossEntropyLoss()
    history = {"train": [], "validation": []}
    for epoch in range(epochs):
        model.train()
        order = torch.randperm(len(X))
        for start in range(0, len(X), batch_size):
            batch = order[start:start + batch_size]
            optimizer.zero_grad()                    # forget the previous gradients
            loss = loss_fn(model(X[batch]), y[batch])
            loss.backward()                          # backpropagation: gradient for every weight
            optimizer.step()                         # move every weight a small step
        history["train"].append(accuracy(model, X, y))
        history["validation"].append(accuracy(model, X_v, y_v))
    return history

torch.manual_seed(42)
mlp_history = train_classifier(mlp, X_fit, y_fit, X_val, y_val)
print("dense network: validation accuracy =", round(mlp_history["validation"][-1], 4))
\end{lstlisting}

\begin{lstlisting}[style=mlpython]
def plot_history(history, title):
    plt.figure(figsize=(5.5, 3.4))
    plt.plot(range(1, len(history["train"]) + 1), history["train"], marker="o", label="training")
    plt.plot(range(1, len(history["validation"]) + 1), history["validation"],
             marker="o", label="validation")
    plt.xlabel("Epoch")
    plt.ylabel("Accuracy")
    plt.title(title)
    plt.legend()
    plt.tight_layout()
    plt.show()

plot_history(mlp_history, "Dense network: accuracy per epoch")
\end{lstlisting}

\textbf{Inspect.} Training accuracy rises while validation plateaus. If validation falls, retain an earlier validated epoch.

\medskip\noindent\textbf{Part C --- A convolutional network.}\par

The CNN applies 16 then 32 filters, pooling after each stage, followed by flattening and dense classification. Print intermediate shapes.

\begin{lstlisting}[style=mlpython]
def make_cnn():
    return nn.Sequential(
        nn.Conv2d(1, 16, kernel_size=3, padding=1), nn.ReLU(), nn.MaxPool2d(2),   # 16 x 4 x 4
        nn.Conv2d(16, 32, kernel_size=3, padding=1), nn.ReLU(), nn.MaxPool2d(2),  # 32 x 2 x 2
        nn.Flatten(),                                                             # 128
        nn.Linear(32 * 2 * 2, 64), nn.ReLU(), nn.Dropout(0.2),
        nn.Linear(64, 10),
    )

cnn = make_cnn()
print("trainable weights:", sum(p.numel() for p in cnn.parameters()))
sample = X_fit[:1]
for layer in cnn:
    sample = layer(sample)
    if isinstance(layer, (nn.Conv2d, nn.MaxPool2d, nn.Flatten, nn.Linear)):
        print(f"{layer.__class__.__name__:10s} -> {tuple(sample.shape[1:])}")
\end{lstlisting}

The same loop trains it; nothing else changes.

\begin{lstlisting}[style=mlpython]
torch.manual_seed(42)
cnn_history = train_classifier(cnn, X_fit, y_fit, X_val, y_val)
print("CNN: validation accuracy =", round(cnn_history["validation"][-1], 4))
plot_history(cnn_history, "CNN: accuracy per epoch")
\end{lstlisting}

Now the official test file, opened once for each model. The third number is the PCA-plus-logistic result from Lab~4 on the same file.

\begin{lstlisting}[style=mlpython]
test_scores = {
    "logistic on PCA (Lab 4)": 0.9477,
    "dense network": accuracy(mlp, X_test, y_test),
    "CNN": accuracy(cnn, X_test, y_test),
}
for name, value in test_scores.items():
    print(f"{name:26s} test accuracy = {value:.4f}")

plt.figure(figsize=(5, 3.2))
plt.bar(list(test_scores), list(test_scores.values()))
plt.ylim(0.9, 1.0)
plt.ylabel("Official test accuracy")
plt.title("Three models, one test file")
plt.tight_layout()
plt.show()
\end{lstlisting}

\textbf{Inspect.} Compare gains across linear, dense, and convolutional models on 1,797 test digits. CNN filters reuse learned patterns across locations.

Inspect the test confusion matrix and misclassified images.

\begin{lstlisting}[style=mlpython]
cnn.eval()
with torch.no_grad():
    test_pred = cnn(X_test).argmax(dim=1)

ConfusionMatrixDisplay.from_predictions(y_test.numpy(), test_pred.numpy(), colorbar=False)
plt.title("CNN confusion matrix on the official test file")
plt.tight_layout()
plt.show()

wrong = torch.nonzero(test_pred != y_test).flatten()[:10]
fig, axes = plt.subplots(1, len(wrong), figsize=(10, 1.6))
for ax, i in zip(axes, wrong):
    ax.imshow(X_test[i, 0], cmap="gray_r")
    ax.set_title(f"{int(y_test[i])} read as {int(test_pred[i])}", fontsize=8)
    ax.axis("off")
plt.suptitle("Ten of the CNN's mistakes", y=1.1)
plt.show()
\end{lstlisting}

\textbf{Inspect.} Check 8/1, 9/5, and 3/8 confusions. Examine images before attributing errors to ambiguity or systematic bias.

Visualize the sixteen first-layer $3\times3$ filters.

\begin{lstlisting}[style=mlpython]
filters = cnn[0].weight.detach()[:, 0]
fig, axes = plt.subplots(2, 8, figsize=(8, 2.4))
for ax, kernel in zip(axes.flat, filters):
    ax.imshow(kernel, cmap="coolwarm", vmin=-filters.abs().max(), vmax=filters.abs().max())
    ax.axis("off")
plt.suptitle("The sixteen 3x3 filters of the first convolutional layer", y=1.02)
plt.show()
\end{lstlisting}

\textbf{Inspect.} Bright/dark filter contrasts can detect oriented edges learned through loss minimization.

\medskip\noindent\textbf{Part D --- Save the weights with a card.}\par

\texttt{torch.save} writes the learned weights (the \emph{state dict}) to a file. Loading them requires building the same architecture first, which is why \texttt{make\_cnn} and the model card belong beside the file. The reloaded model is checked against the in-memory one before the file is trusted.

\begin{lstlisting}[style=mlpython]
weights_path = MODELS / "digits_cnn.pt"
torch.save(cnn.state_dict(), weights_path)

card = {
    "task": "classify 8x8 handwritten digits (0-9)",
    "input": "tensor of shape (n, 1, 8, 8), pixel values divided by 16",
    "architecture": "make_cnn() in Lab 7: conv16-pool-conv32-pool-dense64-dense10",
    "training": {"epochs": 20, "learning_rate": 3e-3, "batch_size": 32, "seed": 42},
    "official_test_accuracy": round(test_scores["CNN"], 4),
    "torch_version": torch.__version__,
}
(MODELS / "digits_cnn_card.json").write_text(json.dumps(card, indent=2))

restored = make_cnn()
restored.load_state_dict(torch.load(weights_path))
restored.eval()
with torch.no_grad():
    same = torch.equal(restored(X_test).argmax(dim=1), test_pred)
size_kb = weights_path.stat().st_size / 1024
print(f"saved {weights_path.name} ({size_kb:.0f} KB); reload reproduces predictions: {same}")
\end{lstlisting}

\begin{tryit}
Compare 5 versus 40 epochs using validation curves. Remove the second convolution/ReLU/pooling block, adjust flattening to $16\times4\times4$, and assess its contribution.
\end{tryit}
\end{labproject}


\labsection{8}{Deep learning II: an LSTM for time series}{sec:lab08-lstm}

\begin{labproject}{Forecast the next hour of bike demand from the previous twenty-four}
\textbf{Goal.} Build chronological windows, compare two naive forecasts with an LSTM, inspect errors, and save weights/scaling.

\medskip\noindent\textbf{Setup.}\par
\begin{lstlisting}[style=mlpython]
import json
import os
from pathlib import Path

# Anaconda's NumPy and pip's PyTorch each ship an OpenMP runtime on Windows; let them coexist.
os.environ.setdefault("KMP_DUPLICATE_LIB_OK", "TRUE")

import matplotlib.pyplot as plt
import numpy as np
import pandas as pd
import torch
from torch import nn

torch.manual_seed(42)
PROJECT_ROOT = Path.cwd().parent
DATA = PROJECT_ROOT / "all_datasets"
MODELS = PROJECT_ROOT / "models"
MODELS.mkdir(exist_ok=True)
\end{lstlisting}

\medskip\noindent\textbf{Part A --- From a series to windows.}\par

Use the 17,379 hourly bike-rental rows. Divide counts by the training-period maximum; future values can exceed the training range.

\begin{lstlisting}[style=mlpython]
hourly = pd.read_csv(DATA / "bike_sharing_dataset_hour.csv")
hourly["dteday"] = pd.to_datetime(hourly["dteday"])
hourly = hourly.sort_values(["dteday", "hr"]).reset_index(drop=True)
series = hourly["cnt"].to_numpy(dtype=np.float32)

n_train_hours = int(len(series) * 0.80)
scale = float(series[:n_train_hours].max())   # a training-set statistic, saved with the model
scaled = series / scale
print("hours:", len(series), " scale factor:", scale)

plt.figure(figsize=(9, 3))
plt.plot(series[:24 * 14])
plt.xlabel("Hour (first two weeks of the series)")
plt.ylabel("Rentals")
plt.title("A daily rhythm the model has to learn")
plt.tight_layout()
plt.show()
\end{lstlisting}

Predict hour $t$ from its preceding 24 values. Split overlapping windows chronologically: first 80\% for training, remainder for validation.

\begin{lstlisting}[style=mlpython]
LOOKBACK = 24
windows = np.stack([scaled[i:i + LOOKBACK] for i in range(len(scaled) - LOOKBACK)])
targets = scaled[LOOKBACK:]
print("windows:", windows.shape, " targets:", targets.shape)

cut = int(len(windows) * 0.80)
X_train = torch.tensor(windows[:cut]).unsqueeze(-1)    # (rows, 24 steps, 1 value per step)
y_train = torch.tensor(targets[:cut])
X_valid = torch.tensor(windows[cut:]).unsqueeze(-1)
y_valid = torch.tensor(targets[cut:])
print("train rows:", len(X_train), " validation rows:", len(X_valid))
print("first window (rentals):", (windows[0] * scale).astype(int)[:8], "...")
print("next hour (the label):", int(targets[0] * scale))
\end{lstlisting}

\medskip\noindent\textbf{Part B --- Two forecasts that cost nothing.}\par

Compare persistence (last value) and same-hour-yesterday (24 hours earlier). Gains over these baselines must justify the network.

\begin{lstlisting}[style=mlpython]
actual = targets[cut:] * scale
persistence = windows[cut:, -1] * scale
yesterday = windows[cut:, 0] * scale

baselines = {
    "persistence (last hour)": np.abs(actual - persistence).mean(),
    "same hour yesterday": np.abs(actual - yesterday).mean(),
}
for name, mae in baselines.items():
    print(f"{name:26s} validation MAE = {mae:.1f} rentals")
\end{lstlisting}

\medskip\noindent\textbf{Part C --- The LSTM.}\par

An LSTM carries 32 hidden values through each 24-step window; a linear layer predicts demand. Reuse Lab~7's loop with squared loss and validation MAE in rentals.

\begin{lstlisting}[style=mlpython]
class LSTMForecaster(nn.Module):
    def __init__(self, hidden_size=32):
        super().__init__()
        self.lstm = nn.LSTM(input_size=1, hidden_size=hidden_size, batch_first=True)
        self.head = nn.Linear(hidden_size, 1)

    def forward(self, x):
        states, _ = self.lstm(x)            # (rows, 24 steps, hidden_size)
        return self.head(states[:, -1, :]).squeeze(-1)   # the final state predicts the next hour

def train_forecaster(model, epochs=8, lr=3e-3, batch_size=64):
    optimizer = torch.optim.Adam(model.parameters(), lr=lr)
    loss_fn = nn.MSELoss()
    history = []
    for epoch in range(epochs):
        model.train()
        order = torch.randperm(len(X_train))
        for start in range(0, len(X_train), batch_size):
            batch = order[start:start + batch_size]
            optimizer.zero_grad()
            loss = loss_fn(model(X_train[batch]), y_train[batch])
            loss.backward()
            optimizer.step()
        model.eval()
        with torch.no_grad():
            mae = (model(X_valid) - y_valid).abs().mean().item() * scale
        history.append(mae)
        print(f"epoch {epoch + 1}: validation MAE = {mae:.1f} rentals")
    return history

torch.manual_seed(42)
model = LSTMForecaster()
print("trainable weights:", sum(p.numel() for p in model.parameters()))
history = train_forecaster(model)
\end{lstlisting}

\begin{lstlisting}[style=mlpython]
plt.figure(figsize=(5.5, 3.4))
plt.plot(range(1, len(history) + 1), history, marker="o", label="LSTM")
for name, mae in baselines.items():
    plt.axhline(mae, linestyle="--", label=name)
plt.xlabel("Epoch")
plt.ylabel("Validation MAE (rentals)")
plt.title("The LSTM has to beat two free forecasts")
plt.legend()
plt.tight_layout()
plt.show()
\end{lstlisting}

\textbf{Inspect.} Validation error falls below both baselines. Use the validation curve to assess whether further epochs help.

Plot forecasts against observations for the first validation week.

\begin{lstlisting}[style=mlpython]
model.eval()
with torch.no_grad():
    forecast = model(X_valid).numpy() * scale

week = slice(0, 24 * 7)
fig, axes = plt.subplots(1, 2, figsize=(11, 3.6), gridspec_kw={"width_ratios": [2, 1]})
axes[0].plot(actual[week], label="actual")
axes[0].plot(forecast[week], label="LSTM forecast")
axes[0].plot(yesterday[week], label="same hour yesterday", alpha=0.5)
axes[0].set_xlabel("Hour of the first validation week")
axes[0].set_ylabel("Rentals")
axes[0].set_title("One-hour-ahead forecasts")
axes[0].legend()
axes[1].scatter(actual, forecast, s=4, alpha=0.3)
axes[1].plot([0, actual.max()], [0, actual.max()], linestyle="--", color="grey")
axes[1].set_xlabel("Actual rentals")
axes[1].set_ylabel("Forecast")
axes[1].set_title("All validation hours")
plt.tight_layout()
plt.show()
\end{lstlisting}

\textbf{Inspect.} Compare rush-hour peaks, nights, and errors at high demand; the largest peaks remain difficult.

\medskip\noindent\textbf{Part D --- Save the weights and the scaling with them.}\par

Save the training maximum and architecture beside the weights. Reload and verify predictions.

\begin{lstlisting}[style=mlpython]
weights_path = MODELS / "bike_lstm.pt"
torch.save(model.state_dict(), weights_path)
card = {
    "task": "forecast next-hour bike rentals from the previous 24 hours",
    "input": "tensor of shape (n, 24, 1): the last 24 hourly counts divided by scale",
    "scale": float(scale),
    "architecture": "LSTMForecaster(hidden_size=32) in Lab 8",
    "training": {"epochs": 8, "learning_rate": 3e-3, "batch_size": 64, "seed": 42},
    "validation_mae_rentals": round(history[-1], 1),
    "baseline_mae_rentals": {k: round(float(v), 1) for k, v in baselines.items()},
    "torch_version": torch.__version__,
}
(MODELS / "bike_lstm_card.json").write_text(json.dumps(card, indent=2))

restored = LSTMForecaster()
restored.load_state_dict(torch.load(weights_path))
restored.eval()
last_day = torch.tensor(scaled[-LOOKBACK:]).reshape(1, LOOKBACK, 1)
with torch.no_grad():
    next_hour = restored(last_day).item() * card["scale"]
print(f"saved {weights_path.name}")
print(f"forecast for the hour after the series ends: {next_hour:.0f} rentals")
\end{lstlisting}

\begin{tryit}
Compare \texttt{LOOKBACK=6} and 48 after eight epochs. Double \texttt{hidden\_size} to 64; compare validation MAE and runtime.
\end{tryit}
\end{labproject}


\begin{quickcheck}
\begin{enumerate}
  \item Why does a network with two linear layers and no activation function have no more power than one linear layer?
  \item What does a convolutional filter share across the image, and why does that reduce the number of weights?
  \item Why must the rows built from a windowed time series be split chronologically rather than at random?
\end{enumerate}
\end{quickcheck}
